\chapter*{Tentang Proyek Logika Terbuka}
\addcontentsline{toc}{chapter}{Tentang Proyek Logika Terbuka}


\textit{Teks Logika Terbuka} adalah buku ajar sumber terbuka dan
kolaboratif tentang metalogika formal dan metode formal, yang dimulai pada
tingkat menengah (yakni, setelah mata kuliah pengantar logika formal).
Walaupun ditujukan kepada pembaca berlatar nonmatematika (khususnya mahasiswa
filsafat dan ilmu komputer), buku ini menyajikan pembahasan yang ketat.

Cakupan beberapa topik yang saat ini disertakan mungkin belum lengkap, dan
banyak bagian masih memerlukan revisi besar. Kami berencana memperluas teks
ini agar kelak mencakup lebih banyak topik. Kami juga berencana menambahkan
sejumlah unsur, seperti glosarium, daftar bacaan lanjutan, catatan historis,
gambar, penjelasan yang lebih baik, bagian yang menjelaskan relevansi
hasil-hasil bagi filsafat, ilmu komputer, dan matematika, serta lebih banyak
soal dan contoh. Jika Anda menemukan kesalahan atau memiliki saran,
\href{https://github.com/OpenLogicProject/OpenLogic/wiki/Contributing}{silakan
beri tahu tim proyek}.

Proyek ini berjalan dalam semangat sumber terbuka. Teks ini tidak hanya
tersedia secara bebas; kami juga menyediakan sumber LaTeX berdasarkan
Lisensi Creative Commons Atribusi, yang memberi setiap orang hak untuk
mengunduh, menggunakan, mengubah, menyusun ulang, mengonversi, dan
mendistribusikan kembali karya kami, asalkan mereka memberikan atribusi yang
sesuai. Informasi tambahan tersedia di situs web Proyek Logika Terbuka,
\href{http://openlogicproject.org/}{openlogicproject.org}.
